% ============================================================================
% CAPÍTULO 8: SISTEMAS METABÓLICOS
% Bioinformática para Biologia de Sistemas
% ============================================================================

% ============================================================================
% TEMPLATE BASE PARA APRESENTAÇÕES EM BEAMER
% Bioinformática para Biologia de Sistemas
% ============================================================================

\documentclass[12pt, aspectratio=169]{beamer}

% ============================================================================
% PACOTES ESSENCIAIS
% ============================================================================
\usepackage[utf8]{inputenc}
\usepackage[T1]{fontenc}
\usepackage[brazilian]{babel}
\usepackage{amsmath, amssymb, amsthm}
\usepackage{graphicx}
\usepackage{xcolor}
\usepackage{tikz}
\usetikzlibrary{arrows.meta, shapes, positioning, calc, decorations.pathreplacing, decorations.shapes, decorations.pathmorphing}
\usepackage{listings}
\usepackage{booktabs}
\usepackage{multirow}
\usepackage{hyperref}
\usepackage{chemfig}
\usepackage{chemarr}

% ============================================================================
% CONFIGURAÇÃO DO TEMA
% ============================================================================
\usetheme{Madrid}
\usecolortheme{default}

% Cores personalizadas
\definecolor{azulescuro}{RGB}{0, 51, 102}
\definecolor{azulclaro}{RGB}{51, 153, 255}
\definecolor{verdeacido}{RGB}{102, 204, 0}
\definecolor{laranjacel}{RGB}{255, 153, 51}
\definecolor{roxodna}{RGB}{153, 51, 255}

\setbeamercolor{structure}{fg=azulescuro}
\setbeamercolor{palette primary}{bg=azulescuro,fg=white}
\setbeamercolor{palette secondary}{bg=azulclaro,fg=white}
\setbeamercolor{palette tertiary}{bg=azulescuro,fg=white}
\setbeamercolor{palette quaternary}{bg=azulclaro,fg=white}
\setbeamercolor{block title}{bg=azulescuro,fg=white}
\setbeamercolor{block body}{bg=azulclaro!10,fg=black}
\setbeamercolor{block title alerted}{bg=red!80,fg=white}
\setbeamercolor{block body alerted}{bg=red!10,fg=black}
\setbeamercolor{block title example}{bg=verdeacido!80,fg=white}
\setbeamercolor{block body example}{bg=verdeacido!10,fg=black}

% ============================================================================
% CONFIGURAÇÃO DE CÓDIGO
% ============================================================================
\lstset{
    language=Python,
    basicstyle=\ttfamily\small,
    keywordstyle=\color{azulescuro}\bfseries,
    commentstyle=\color{verdeacido},
    stringstyle=\color{laranjacel},
    numbers=left,
    numberstyle=\tiny\color{gray},
    stepnumber=1,
    numbersep=5pt,
    backgroundcolor=\color{white},
    showspaces=false,
    showstringspaces=false,
    showtabs=false,
    frame=single,
    rulecolor=\color{black},
    tabsize=2,
    captionpos=b,
    breaklines=true,
    breakatwhitespace=false,
    escapeinside={\%*}{*)},
    xleftmargin=10pt,
    xrightmargin=10pt
}

% Definir estilo para R
\lstdefinestyle{Rstyle}{
    language=R,
    basicstyle=\ttfamily\small,
    keywordstyle=\color{azulescuro}\bfseries,
    commentstyle=\color{verdeacido},
    stringstyle=\color{laranjacel}
}

% Definir estilo para MATLAB
\lstdefinestyle{MATLABstyle}{
    language=Matlab,
    basicstyle=\ttfamily\small,
    keywordstyle=\color{azulescuro}\bfseries,
    commentstyle=\color{verdeacido},
    stringstyle=\color{laranjacel}
}

% ============================================================================
% COMANDOS PERSONALIZADOS
% ============================================================================

% Comando para equações destacadas
\newcommand{\destaque}[1]{%
    \begin{alertblock}{}
        \begin{center}
            #1
        \end{center}
    \end{alertblock}
}

% Comando para definições
\newcommand{\definicao}[2]{%
    \begin{block}{Definição: #1}
        #2
    \end{block}
}

% Comando para exemplos
\newcommand{\exemplo}[2]{%
    \begin{exampleblock}{Exemplo: #1}
        #2
    \end{exampleblock}
}

% Comando para observações importantes
\newcommand{\importante}[1]{%
    \begin{alertblock}{Importante!}
        #1
    \end{alertblock}
}

% Comandos matemáticos úteis
\newcommand{\R}{\mathbb{R}}
\newcommand{\N}{\mathbb{N}}
\newcommand{\Z}{\mathbb{Z}}
\newcommand{\dif}{\mathrm{d}}
\newcommand{\parcial}[2]{\frac{\partial #1}{\partial #2}}
\newcommand{\derivada}[2]{\frac{\dif #1}{\dif #2}}
\newcommand{\vect}[1]{\boldsymbol{#1}}

% Comandos biológicos
\newcommand{\gene}[1]{\textit{#1}}
\newcommand{\proteina}[1]{\textsc{#1}}
\newcommand{\especie}[1]{\textit{#1}}

% ============================================================================
% CONFIGURAÇÃO DE RODAPÉ
% ============================================================================
\setbeamertemplate{footline}{
  \leavevmode%
  \hbox{%
  \begin{beamercolorbox}[wd=.33\paperwidth,ht=2.25ex,dp=1ex,center]{author in head/foot}%
    \usebeamerfont{author in head/foot}\insertshortauthor
  \end{beamercolorbox}%
  \begin{beamercolorbox}[wd=.34\paperwidth,ht=2.25ex,dp=1ex,center]{title in head/foot}%
    \usebeamerfont{title in head/foot}\insertshorttitle
  \end{beamercolorbox}%
  \begin{beamercolorbox}[wd=.33\paperwidth,ht=2.25ex,dp=1ex,right]{date in head/foot}%
    \usebeamerfont{date in head/foot}\insertshortdate{}\hspace*{2em}
    \insertframenumber{} / \inserttotalframenumber\hspace*{2ex}
  \end{beamercolorbox}}%
  \vskip0pt%
}

% ============================================================================
% CONFIGURAÇÃO DE NAVEGAÇÃO
% ============================================================================
\setbeamertemplate{navigation symbols}{}

% ============================================================================
% CONFIGURAÇÃO DE ITEMIZE/ENUMERATE
% ============================================================================
\setbeamertemplate{itemize items}[circle]
\setbeamertemplate{enumerate items}[default]

% ============================================================================
% FIM DO TEMPLATE
% ============================================================================


% ============================================================================
% INFORMAÇÕES DO DOCUMENTO
% ============================================================================
\title[Cap. 8: Sistemas Metabólicos]{Capítulo 8: Sistemas Metabólicos}
\subtitle{Bioinformática para Biologia de Sistemas}
\author{Ney Lemke}
\institute{DBF-Unesp}
\date{\today}

\begin{document}

% ============================================================================
% SLIDE DE TÍTULO
% ============================================================================
\begin{frame}
\titlepage
\end{frame}

% ============================================================================
% SUMÁRIO
% ============================================================================
\begin{frame}{Sumário}
\tableofcontents
\end{frame}

% ============================================================================
% SEÇÃO 1: INTRODUÇÃO
% ============================================================================
\section{Introdução}

\begin{frame}{Visão Geral do Capítulo}
\begin{block}{Objetivos de Aprendizagem}
\begin{itemize}
    \item Compreender a estrutura e dinâmica de redes metabólicas
    \item Modelar cinéticas de reações enzimáticas
    \item Aplicar análise de fluxo metabólico (MFA)
    \item Utilizar análise de balanço de fluxo (FBA)
    \item Interpretar dados metabolômicos
\end{itemize}
\end{block}

\vspace{0.3cm}

\begin{exampleblock}{Por que Estudar Metabolismo?}
O metabolismo é fundamental para:
\begin{itemize}
    \item Produção de energia e biomoléculas
    \item Engenharia metabólica e biotecnologia
    \item Compreensão de doenças metabólicas
    \item Design racional de terapias
\end{itemize}
\end{exampleblock}
\end{frame}

\begin{frame}{O Papel Central do Metabolismo}
\definicao{Metabolismo}{
Conjunto de reações químicas interconectadas que convertem nutrientes em energia e blocos construtores para crescimento e manutenção celular.
}

\vspace{0.3cm}

\begin{center}
\begin{tikzpicture}[node distance=2.5cm, scale=0.9, transform shape]
    \node[draw, circle, fill=laranjacel!30, text width=2cm, align=center] (nut) {Nutrientes};
    \node[draw, circle, fill=azulclaro!30, text width=2cm, align=center, right=of nut] (met) {Metabolismo};
    \node[draw, circle, fill=verdeacido!30, text width=1.5cm, align=center, above right=1cm and 1.5cm of met] (atp) {Energia\\(ATP)};
    \node[draw, circle, fill=roxodna!30, text width=1.5cm, align=center, below right=1cm and 1.5cm of met] (bio) {Biomassa};

    \draw[->, thick, azulescuro, line width=1.5pt] (nut) -- (met);
    \draw[->, thick, azulescuro, line width=1.5pt] (met) -- (atp);
    \draw[->, thick, azulescuro, line width=1.5pt] (met) -- (bio);
\end{tikzpicture}
\end{center}

\importante{
Redes metabólicas são sistemas complexos com milhares de reações interconectadas
}
\end{frame}

\begin{frame}{Redes Metabólicas como Sistemas Complexos}
\begin{columns}
\column{0.5\textwidth}
\textbf{Características:}
\begin{itemize}
    \item \textbf{Escala:} centenas a milhares de reações
    \item \textbf{Conectividade:} altamente interconectadas
    \item \textbf{Regulação:} múltiplos níveis de controle
    \item \textbf{Robustez:} redundância e flexibilidade
    \item \textbf{Otimalidade:} eficiência evolutiva
\end{itemize}

\column{0.5\textwidth}
\textbf{Níveis de Análise:}
\begin{itemize}
    \item \textbf{Cinética:} dinâmica de reações individuais
    \item \textbf{Estequiométrica:} balanços de massa
    \item \textbf{Topológica:} estrutura da rede
    \item \textbf{Regulatória:} controle de fluxo
    \item \textbf{Evolutiva:} adaptação e otimização
\end{itemize}
\end{columns}

\vspace{0.3cm}

\begin{exampleblock}{Abordagens Computacionais}
Modelagem cinética, FBA, MFA, análise de modos elementares, metabolômica
\end{exampleblock}
\end{frame}

% ============================================================================
% SEÇÃO 2: CINÉTICAS DE REAÇÃO METABÓLICA
% ============================================================================
\section{Cinéticas de Reação Metabólica}

\begin{frame}{Visão Geral das Cinéticas Enzimáticas}
\begin{block}{Por que Modelar Cinéticas?}
\begin{itemize}
    \item Predizer comportamento dinâmico de vias metabólicas
    \item Identificar passos limitantes
    \item Avaliar estratégias de engenharia metabólica
    \item Compreender regulação alostérica
\end{itemize}
\end{block}

\vspace{0.3cm}

\begin{columns}
\column{0.5\textwidth}
\textbf{Tipos de Modelos:}
\begin{itemize}
    \item Ação de massas
    \item Michaelis-Menten
    \item Hill cooperativo
    \item Convenience kinetics
    \item GMA/S-systems
\end{itemize}

\column{0.5\textwidth}
\begin{center}
\begin{tikzpicture}[scale=0.7]
    \draw[->] (0,0) -- (5,0) node[right] {$[S]$};
    \draw[->] (0,0) -- (0,4) node[above] {$v$};

    % Michaelis-Menten
    \draw[thick, azulescuro, domain=0:4.5, smooth, samples=100]
        plot (\x, {3*\x/(1+\x)});
    \draw[dashed, verdeacido] (0,3) -- (5,3) node[right, font=\tiny] {$V_{max}$};
    \draw[dashed, laranjacel] (1,0) -- (1,1.5) node[right, font=\tiny] {$K_M$};

    \node[below] at (2.5,-0.3) {\small Michaelis-Menten};
\end{tikzpicture}
\end{center}
\end{columns}
\end{frame}

\begin{frame}{Cinética de Ação de Massas}
\begin{block}{Lei de Ação de Massas}
Para uma reação elementar $A + B \xrightarrow{k} C$:
\begin{equation*}
v = k[A][B]
\end{equation*}
\end{block}

\textbf{Características:}
\begin{itemize}
    \item Aplica-se a reações elementares (não-catalisadas)
    \item Linear em cada reagente
    \item Base para derivar outras cinéticas
\end{itemize}

\vspace{0.3cm}

\exemplo{Sistema de Duas Reações}{
\begin{align*}
A &\xrightarrow{k_1} B \quad v_1 = k_1[A]\\
B &\xrightarrow{k_2} C \quad v_2 = k_2[B]
\end{align*}

Sistema de EDOs:
\begin{align*}
\derivada{[A]}{t} &= -k_1[A]\\
\derivada{[B]}{t} &= k_1[A] - k_2[B]\\
\derivada{[C]}{t} &= k_2[B]
\end{align*}
}
\end{frame}

\begin{frame}{Cinética de Michaelis-Menten}
\definicao{Michaelis-Menten}{
Modelo padrão para catálise enzimática com um substrato:
\begin{equation*}
v = \frac{V_{max}[S]}{K_M + [S]}
\end{equation*}
onde $V_{max} = k_{cat}[E]_T$ e $K_M = \frac{k_{-1} + k_{cat}}{k_1}$
}

\vspace{0.3cm}

\begin{columns}
\column{0.5\textwidth}
\textbf{Derivação:}
\begin{equation*}
E + S \xrightleftharpoons[k_{-1}]{k_1} ES \xrightarrow{k_{cat}} E + P
\end{equation*}

Hipótese de quasi-estado estacionário:
\begin{equation*}
\derivada{[ES]}{t} \approx 0
\end{equation*}

\column{0.5\textwidth}
\textbf{Parâmetros:}
\begin{itemize}
    \item $K_M$: afinidade (baixo $K_M$ = alta afinidade)
    \item $V_{max}$: capacidade catalítica máxima
    \item $k_{cat}$: número de turnover
    \item $k_{cat}/K_M$: eficiência catalítica
\end{itemize}
\end{columns}
\end{frame}

\begin{frame}{Equação de Hill: Cooperatividade}
\begin{block}{Cinética de Hill}
Para enzimas com múltiplos sítios de ligação e cooperatividade:
\begin{equation*}
v = \frac{V_{max}[S]^n}{K_0.5^n + [S]^n}
\end{equation*}
onde $n$ é o coeficiente de Hill e $K_{0.5}$ é a constante de meia-saturação.
\end{block}

\vspace{0.3cm}

\begin{columns}
\column{0.5\textwidth}
\textbf{Interpretação de $n$:}
\begin{itemize}
    \item $n = 1$: sem cooperatividade (MM)
    \item $n > 1$: cooperatividade positiva
    \item $n < 1$: cooperatividade negativa
\end{itemize}

\exemplo{Hemoglobina}{
Ligação de O$_2$ tem $n \approx 2.8$ (cooperatividade positiva forte)
}

\column{0.5\textwidth}
\begin{center}
\begin{tikzpicture}[scale=0.65]
    \draw[->] (0,0) -- (5,0) node[right] {$[S]$};
    \draw[->] (0,0) -- (0,4) node[above] {$v$};

    % n=1 (MM)
    \draw[thick, azulescuro, domain=0:4.5, smooth, samples=100]
        plot (\x, {3*\x/(1+\x)}) node[right, font=\tiny] {$n=1$};

    % n=3
    \draw[thick, verdeacido, domain=0:4.5, smooth, samples=100]
        plot (\x, {3*\x^3/(1+\x^3)}) node[right, font=\tiny] {$n=3$};

    % n=0.5
    \draw[thick, laranjacel, domain=0:4.5, smooth, samples=100]
        plot (\x, {3*sqrt(\x)/(1+sqrt(\x))}) node[right, font=\tiny] {$n=0.5$};

    \node[below] at (2.5,-0.3) {\small Comparação};
\end{tikzpicture}
\end{center}
\end{columns}
\end{frame}

\begin{frame}{Michaelis-Menten Reversível}
\begin{block}{Reação Bidirecional}
Para reação reversível $S \rightleftharpoons P$:
\destaque{
\begin{equation*}
v = \frac{V_f\frac{[S]}{K_{MS}} - V_r\frac{[P]}{K_{MP}}}{1 + \frac{[S]}{K_{MS}} + \frac{[P]}{K_{MP}}}
\end{equation*}
}

onde:
\begin{itemize}
    \item $V_f$, $V_r$: velocidades máximas forward e reverse
    \item $K_{MS}$, $K_{MP}$: constantes de Michaelis para S e P
\end{itemize}
\end{block}

\begin{exampleblock}{Relação com Termodinâmica}
No equilíbrio ($v = 0$):
\begin{equation*}
K_{eq} = \frac{[P]_{eq}}{[S]_{eq}} = \frac{V_f K_{MP}}{V_r K_{MS}}
\end{equation*}
\end{exampleblock}
\end{frame}

\begin{frame}{Convenience Kinetics}
\begin{block}{Forma Geral}
Aproximação prática que incorpora inibição por produto e regulação:
\begin{equation*}
v = E_T \cdot k_{cat} \cdot \frac{\frac{[S]}{K_S}\left(1 - \frac{[P]}{[S]K_{eq}}\right)}{\left(1 + \frac{[S]}{K_S} + \frac{[P]}{K_P}\right)\left(1 + \sum_i \frac{[I_i]}{K_{Ii}}\right)}
\end{equation*}
\end{block}

\textbf{Vantagens:}
\begin{itemize}
    \item Termodinamicamente consistente
    \item Incorpora múltiplos efetores
    \item Parametrização mais simples que modelos detalhados
    \item Amplamente usada em modelos de larga escala
\end{itemize}

\textbf{Desvantagens:}
\begin{itemize}
    \item Menos mecanística
    \item Requer ajuste cuidadoso de parâmetros
\end{itemize}
\end{frame}

\begin{frame}{Generalized Mass Action (GMA) e S-Systems}
\begin{columns}
\column{0.5\textwidth}
\begin{block}{GMA (Lei de Potência)}
\begin{equation*}
v_i = \alpha_i \prod_{j=1}^{n} X_j^{g_{ij}}
\end{equation*}

\textbf{Sistema GMA:}
\begin{equation*}
\derivada{X_i}{t} = \sum_k \alpha_k \prod_j X_j^{g_{kj}} - \sum_l \beta_l \prod_j X_j^{h_{lj}}
\end{equation*}
\end{block}

\column{0.5\textwidth}
\begin{block}{S-Systems}
Forma especial de GMA:
\begin{equation*}
\derivada{X_i}{t} = \alpha_i \prod_{j=1}^{n} X_j^{g_{ij}} - \beta_i \prod_{j=1}^{n} X_j^{h_{ij}}
\end{equation*}

Separação em:
\begin{itemize}
    \item Termo de produção agregado
    \item Termo de degradação agregado
\end{itemize}
\end{block}
\end{columns}

\vspace{0.3cm}

\importante{
S-systems permitem análise algébrica de estado estacionário e são base da Biochemical Systems Theory (BST)
}
\end{frame}

\begin{frame}{Comparação de Formalismos Cinéticos}
\begin{center}
\begin{tabular}{lccc}
\toprule
\textbf{Tipo} & \textbf{Complexidade} & \textbf{Mecanística} & \textbf{Parametrização} \\
\midrule
Ação de massas & Baixa & Alta & Fácil \\
Michaelis-Menten & Média & Alta & Média \\
Hill & Média & Média & Média \\
Convenience & Alta & Média & Difícil \\
GMA/S-systems & Média & Baixa & Difícil \\
\bottomrule
\end{tabular}
\end{center}

\vspace{0.3cm}

\begin{exampleblock}{Escolha do Modelo}
\begin{itemize}
    \item \textbf{Pequenas vias detalhadas}: Michaelis-Menten, Hill
    \item \textbf{Vias de médio porte}: Convenience kinetics
    \item \textbf{Análise de estado estacionário}: S-systems, FBA
    \item \textbf{Modelos de genoma completo}: FBA (sem cinética)
\end{itemize}
\end{exampleblock}
\end{frame}

\begin{frame}[fragile]{Implementação: Cinética de Michaelis-Menten}
\begin{lstlisting}[language=Python, caption=Modelo cinético simples]
import numpy as np
import matplotlib.pyplot as plt
from scipy.integrate import odeint

def michaelis_menten(S, t, Vmax, Km, k_deg):
    """
    Modelo simples: S -> P com MM kinetics
    S: substrato, P: produto
    """
    v = Vmax * S / (Km + S)
    dSdt = -v - k_deg * S  # consumo + degradacao
    return dSdt

# Parametros
Vmax = 10.0   # umol/min
Km = 2.0      # mM
k_deg = 0.1   # 1/min

# Condicao inicial
S0 = 20.0     # mM

# Tempo
t = np.linspace(0, 50, 500)

# Resolver
S = odeint(michaelis_menten, S0, t, args=(Vmax, Km, k_deg))

# Plotar
plt.figure(figsize=(10, 6))
plt.plot(t, S, 'b-', linewidth=2, label='[S](t)')
plt.xlabel('Tempo (min)', fontsize=12)
plt.ylabel('Concentracao [S] (mM)', fontsize=12)
plt.title('Cinetica de Michaelis-Menten', fontsize=14)
plt.legend()
plt.grid(True, alpha=0.3)
\end{lstlisting}
\end{frame}

\begin{frame}[fragile]{Implementação: Via Metabólica Simples}
\begin{lstlisting}[language=Python, caption=Via de três passos, basicstyle=\ttfamily\footnotesize]
def metabolic_pathway(y, t, params):
    """
    Via metabolica: A -> B -> C -> D
    y = [A, B, C, D]
    """
    A, B, C, D = y
    Vmax1, Km1, Vmax2, Km2, Vmax3, Km3 = params

    # Reacoes com MM kinetics
    v1 = Vmax1 * A / (Km1 + A)
    v2 = Vmax2 * B / (Km2 + B)
    v3 = Vmax3 * C / (Km3 + C)

    # EDOs
    dAdt = -v1
    dBdt = v1 - v2
    dCdt = v2 - v3
    dDdt = v3

    return [dAdt, dBdt, dCdt, dDdt]

# Parametros [Vmax1, Km1, Vmax2, Km2, Vmax3, Km3]
params = [10.0, 1.0, 8.0, 1.5, 5.0, 2.0]

# Condicoes iniciais [A0, B0, C0, D0]
y0 = [50.0, 0.0, 0.0, 0.0]

# Simular
t = np.linspace(0, 20, 500)
sol = odeint(metabolic_pathway, y0, t, args=(params,))

# Plotar todas as especies
plt.plot(t, sol[:, 0], 'b-', label='A', linewidth=2)
plt.plot(t, sol[:, 1], 'r-', label='B', linewidth=2)
plt.plot(t, sol[:, 2], 'g-', label='C', linewidth=2)
plt.plot(t, sol[:, 3], 'm-', label='D', linewidth=2)
plt.legend()
\end{lstlisting}
\end{frame}

% ============================================================================
% SEÇÃO 3: ANÁLISE DE FLUXO METABÓLICO (MFA)
% ============================================================================
\section{Análise de Fluxo Metabólico (MFA)}

\begin{frame}{Análise de Fluxo Metabólico (MFA)}
\definicao{Metabolic Flux Analysis (MFA)}{
Técnica experimental e computacional para determinar taxas de reação (fluxos) em uma rede metabólica usando dados de medições e balanços de massa.
}

\vspace{0.3cm}

\begin{block}{Princípios Fundamentais}
\begin{itemize}
    \item \textbf{Conservação de massa:} entrada = saída + acúmulo
    \item \textbf{Estado pseudo-estacionário:} metabólitos intracelulares constantes
    \item \textbf{Medições externas:} taxas de consumo/produção
    \item \textbf{Marcação isotópica:} rastreamento de carbono-13
\end{itemize}
\end{block}

\begin{exampleblock}{Aplicações}
Engenharia metabólica, identificação de gargalos, otimização de bioprocessos
\end{exampleblock}
\end{frame}

\begin{frame}{Representação Estequiométrica}
\begin{block}{Matriz Estequiométrica $\mathbf{S}$}
Representa a topologia da rede:
\begin{itemize}
    \item Linhas: metabólitos ($m$)
    \item Colunas: reações ($n$)
    \item $S_{ij}$: coeficiente estequiométrico do metabólito $i$ na reação $j$
\end{itemize}
\end{block}

\exemplo{Rede Simples}{
Reações:
\begin{align*}
v_1&: \quad A \rightarrow B\\
v_2&: \quad B \rightarrow C\\
v_3&: \quad B \rightarrow D
\end{align*}

Matriz estequiométrica:
\begin{equation*}
\mathbf{S} = \begin{bmatrix}
-1 & 0 & 0 \\
1 & -1 & -1 \\
0 & 1 & 0 \\
0 & 0 & 1
\end{bmatrix}
\begin{matrix}
\leftarrow A\\
\leftarrow B\\
\leftarrow C\\
\leftarrow D
\end{matrix}
\end{equation*}
}
\end{frame}

\begin{frame}{Balanço de Massa}
\begin{block}{Equação de Balanço}
Para cada metabólito intracelular:
\destaque{
\begin{equation*}
\derivada{\vect{c}}{t} = \mathbf{S} \cdot \vect{v}
\end{equation*}
}

onde:
\begin{itemize}
    \item $\vect{c}$: vetor de concentrações ($m \times 1$)
    \item $\mathbf{S}$: matriz estequiométrica ($m \times n$)
    \item $\vect{v}$: vetor de fluxos ($n \times 1$)
\end{itemize}
\end{block}

\begin{block}{Hipótese de Estado Estacionário}
Assumindo $\derivada{\vect{c}}{t} \approx \vect{0}$:
\begin{equation*}
\mathbf{S} \cdot \vect{v} = \vect{0}
\end{equation*}

Sistema de equações algébricas lineares!
\end{block}
\end{frame}

\begin{frame}{Determinação de Fluxos}
\begin{block}{Sistema Subdeterminado}
Tipicamente, $n > m$ (mais fluxos que metabólitos):
\begin{itemize}
    \item Sistema com infinitas soluções
    \item Espaço nulo de $\mathbf{S}$ não-trivial
    \item Requer medições adicionais
\end{itemize}
\end{block}

\vspace{0.3cm}

\textbf{Estratégias:}
\begin{enumerate}
    \item \textbf{Medições diretas:} taxas de consumo/produção extracelular
    \item \textbf{Marcação isotópica:} $^{13}$C-MFA
    \item \textbf{Restrições adicionais:} razões de fluxo, irreversibilidade
    \item \textbf{Otimização:} minimizar função objetivo
\end{enumerate}

\importante{
$^{13}$C-MFA usa padrões de marcação para resolver ambiguidades
}
\end{frame}

\begin{frame}{$^{13}$C-MFA: Conceito}
\begin{block}{Metabolic Flux Analysis com $^{13}$C}
\begin{itemize}
    \item Alimentar células com substrato marcado (e.g., $[1,2\text{-}^{13}$C]glicose)
    \item Medir distribuição de marcação em metabólitos via MS ou NMR
    \item Simular propagação de marcação pela rede
    \item Ajustar fluxos para reproduzir padrões experimentais
\end{itemize}
\end{block}

\vspace{0.3cm}

\begin{columns}
\column{0.5\textwidth}
\textbf{Informação Adicional:}
\begin{itemize}
    \item Direção de fluxos reversíveis
    \item Fluxos em vias paralelas
    \item Ciclos e recirculação
\end{itemize}

\column{0.5\textwidth}
\begin{center}
\begin{tikzpicture}[scale=0.7, every node/.style={font=\tiny}]
    \node[draw, circle, fill=verdeacido!30] (glc) at (0,0) {$^{13}$C-Glc};
    \node[draw, circle, fill=azulclaro!20] (g6p) at (2,0) {G6P};
    \node[draw, circle, fill=azulclaro!20] (pyr) at (2,-2) {Pyr};
    \node[draw, circle, fill=azulclaro!20] (accoa) at (4,-2) {AcCoA};

    \draw[->, thick] (glc) -- (g6p);
    \draw[->, thick] (g6p) -- (pyr);
    \draw[->, thick] (pyr) -- (accoa);

    \node[right, font=\scriptsize] at (4.5,-1) {$^{13}$C label\\tracking};
\end{tikzpicture}
\end{center}
\end{columns}
\end{frame}

\begin{frame}{Metodologia de $^{13}$C-MFA}
\begin{enumerate}
    \item \textbf{Design experimental}
    \begin{itemize}
        \item Escolher substrato marcado apropriado
        \item Cultivar células em estado estacionário
        \item Medir taxas extracelulares e padrões de marcação
    \end{itemize}

    \item \textbf{Modelo de rede}
    \begin{itemize}
        \item Definir topologia metabólica
        \item Mapear transições atômicas
        \item Estabelecer equações de balanço
    \end{itemize}

    \item \textbf{Simulação de marcação}
    \begin{itemize}
        \item Calcular distribuição de marcação para fluxos dados
        \item Usar modelos de isotopômeros ou cumômeros
    \end{itemize}

    \item \textbf{Estimação de fluxos}
    \begin{itemize}
        \item Minimizar desvio entre dados simulados e experimentais
        \item Usar mínimos quadrados ponderados
        \item Análise de incerteza estatística
    \end{itemize}
\end{enumerate}
\end{frame}

\begin{frame}{Exemplo de MFA: Rede Simples}
\begin{columns}
\column{0.5\textwidth}
\textbf{Rede:}
\begin{align*}
v_1&: A_{ext} \rightarrow A\\
v_2&: A \rightarrow B\\
v_3&: B \rightarrow C\\
v_4&: B \rightarrow D\\
v_5&: C \rightarrow C_{ext}\\
v_6&: D \rightarrow D_{ext}
\end{align*}

\textbf{Balanços (steady-state):}
\begin{align*}
A&: v_1 - v_2 = 0\\
B&: v_2 - v_3 - v_4 = 0\\
C&: v_3 - v_5 = 0\\
D&: v_4 - v_6 = 0
\end{align*}

\column{0.5\textwidth}
\textbf{Medições:}
\begin{itemize}
    \item $v_1 = 10$ (uptake de A)
    \item $v_5 = 6$ (produção de C)
    \item $v_6 = 4$ (produção de D)
\end{itemize}

\textbf{Solução:}
\begin{align*}
v_1 &= 10\\
v_2 &= 10\\
v_3 &= 6\\
v_4 &= 4\\
v_5 &= 6\\
v_6 &= 4
\end{align*}

Sistema completamente determinado!
\end{columns}
\end{frame}

\begin{frame}[fragile]{Implementação: MFA Básico}
\begin{lstlisting}[language=Python, caption=Resolução de fluxos por balanço de massa]
import numpy as np
from scipy.linalg import lstsq

# Matriz estequiometrica
# Metabolitos: A, B, C, D (internos)
# Reacoes: v1, v2, v3, v4, v5, v6
S = np.array([
    [ 1, -1,  0,  0,  0,  0],  # A
    [ 0,  1, -1, -1,  0,  0],  # B
    [ 0,  0,  1,  0, -1,  0],  # C
    [ 0,  0,  0,  1,  0, -1]   # D
])

# Medicoes: v1=10, v5=6, v6=4
# Reordenar: fluxos medidos vs nao-medidos
# Fluxos nao-medidos: v2, v3, v4
# S*v = 0 => S_unmeas*v_unmeas = -S_meas*v_meas

S_meas = S[:, [0, 4, 5]]      # colunas v1, v5, v6
S_unmeas = S[:, [1, 2, 3]]    # colunas v2, v3, v4
v_meas = np.array([10, 6, 4])

# Resolver S_unmeas * v_unmeas = -S_meas * v_meas
b = -S_meas @ v_meas
v_unmeas, residuals, rank, s = lstsq(S_unmeas, b)

# Reconstruir vetor completo
v_all = np.zeros(6)
v_all[[0, 4, 5]] = v_meas
v_all[[1, 2, 3]] = v_unmeas

print("Fluxos calculados:")
for i, flux in enumerate(v_all, 1):
    print(f"v{i} = {flux:.2f}")
\end{lstlisting}
\end{frame}

% ============================================================================
% SEÇÃO 4: ANÁLISE DE BALANÇO DE FLUXO (FBA)
% ============================================================================
\section{Análise de Balanço de Fluxo (FBA)}

\begin{frame}{Análise de Balanço de Fluxo (FBA)}
\definicao{Flux Balance Analysis (FBA)}{
Método computacional baseado em otimização linear para predizer fluxos metabólicos em redes de escala genômica, sem requerer parâmetros cinéticos.
}

\vspace{0.3cm}

\begin{block}{Características}
\begin{itemize}
    \item \textbf{Baseada em restrições:} estequiometria + limites de fluxo
    \item \textbf{Estado estacionário:} $\mathbf{S} \cdot \vect{v} = \vect{0}$
    \item \textbf{Função objetivo:} tipicamente maximizar biomassa
    \item \textbf{Escala genômica:} milhares de reações
    \item \textbf{Computacionalmente eficiente:} programação linear
\end{itemize}
\end{block}

\importante{
FBA permite análise de metabolismo em escala de sistemas sem conhecer todas as constantes cinéticas
}
\end{frame}

\begin{frame}{Modelos Metabólicos de Escala Genômica (GEMs)}
\begin{block}{Genome-scale Metabolic Models}
\begin{itemize}
    \item Reconstruídos a partir de anotação genômica
    \item Integram conhecimento bioquímico e fisiológico
    \item Incluem todas as reações conhecidas de um organismo
    \item Validados contra dados experimentais
\end{itemize}
\end{block}

\begin{exampleblock}{Exemplos de GEMs}
\begin{itemize}
    \item \textbf{E. coli:} iJO1366 (2583 reações, 1805 genes)
    \item \textbf{S. cerevisiae:} Yeast8 (4059 reações, 1150 genes)
    \item \textbf{H. sapiens:} Recon3D (13543 reações, 3288 genes)
\end{itemize}
\end{exampleblock}

\vspace{0.2cm}

\begin{center}
\small
Disponíveis em: BiGG Models (\url{http://bigg.ucsd.edu})
\end{center}
\end{frame}

\begin{frame}{Formulação Matemática de FBA}
\begin{block}{Problema de Otimização Linear}
\destaque{
\begin{align*}
\text{Maximizar:} \quad & \vect{c}^T \vect{v}\\
\text{Sujeito a:} \quad & \mathbf{S} \cdot \vect{v} = \vect{0}\\
& \vect{v}_{min} \leq \vect{v} \leq \vect{v}_{max}
\end{align*}
}

onde:
\begin{itemize}
    \item $\vect{c}$: vetor de coeficientes da função objetivo
    \item $\vect{v}$: vetor de fluxos (variáveis)
    \item $\mathbf{S}$: matriz estequiométrica
    \item $\vect{v}_{min}$, $\vect{v}_{max}$: limites de fluxo
\end{itemize}
\end{block}

\begin{exampleblock}{Função Objetivo Típica}
Maximizar crescimento: $c_i = 1$ para reação de biomassa, $c_i = 0$ para outras
\end{exampleblock}
\end{frame}

\begin{frame}{Reação de Biomassa}
\begin{block}{Pseudo-reação de Crescimento}
Representa a composição biomolecular da célula:
\begin{equation*}
\sum_i a_i \cdot M_i \rightarrow 1 \text{ g biomassa}
\end{equation*}

onde $M_i$ são precursores (aminoácidos, nucleotídeos, lipídios, etc.)
\end{block}

\exemplo{Componentes Típicos}{
\small
\begin{itemize}
    \item Proteínas: $\sim$55\%
    \item RNA: $\sim$20\%
    \item DNA: $\sim$3\%
    \item Lipídios: $\sim$9\%
    \item Carboidratos: $\sim$3\%
    \item Pequenas moléculas, íons: $\sim$10\%
\end{itemize}
}

\importante{
A maximização de biomassa é apoiada por evidências evolutivas e valida bem contra dados experimentais
}
\end{frame}

\begin{frame}{Restrições e Limites de Fluxo}
\begin{block}{Tipos de Restrições}
\begin{enumerate}
    \item \textbf{Termodinâmicas:} irreversibilidade
    \begin{itemize}
        \item Reações irreversíveis: $v_i \geq 0$
        \item Reações reversíveis: $v_{min} \leq v_i \leq v_{max}$
    \end{itemize}

    \item \textbf{Capacidade enzimática:} $v_i \leq v_{i,max}$

    \item \textbf{Uptake de nutrientes:} taxas medidas experimentalmente
    \begin{itemize}
        \item Exemplo: $v_{glucose} = -10$ mmol/gDW/h
    \end{itemize}

    \item \textbf{ATP de manutenção:} consumo basal de energia
    \begin{itemize}
        \item Exemplo: $v_{ATPM} \geq 8.39$ mmol/gDW/h
    \end{itemize}
\end{enumerate}
\end{block}

\textbf{Limites padrão:}
\begin{itemize}
    \item Irreversíveis: $0 \leq v_i \leq 1000$
    \item Reversíveis: $-1000 \leq v_i \leq 1000$
\end{itemize}
\end{frame}

\begin{frame}{Análise de Variabilidade de Fluxo (FVA)}
\definicao{Flux Variability Analysis}{
Determina o intervalo de valores possíveis para cada fluxo enquanto mantém a função objetivo em seu valor ótimo (ou próximo).
}

\vspace{0.3cm}

\begin{block}{Formulação}
Para cada reação $j$:
\begin{align*}
\text{Minimizar/Maximizar:} \quad & v_j\\
\text{Sujeito a:} \quad & \mathbf{S} \cdot \vect{v} = \vect{0}\\
& \vect{v}_{min} \leq \vect{v} \leq \vect{v}_{max}\\
& \vect{c}^T \vect{v} \geq \gamma \cdot Z_{opt}
\end{align*}

onde $Z_{opt}$ é o valor objetivo ótimo e $\gamma \in [0,1]$ (tipicamente 0.9-1.0)
\end{block}

\begin{exampleblock}{Interpretação}
\begin{itemize}
    \item Intervalo estreito: fluxo bem determinado
    \item Intervalo amplo: múltiplas soluções alternativas
\end{itemize}
\end{exampleblock}
\end{frame}

\begin{frame}{Predição de Essencialidade Gênica}
\begin{block}{Análise de Gene Knockout}
Simular deleção de gene removendo reações associadas:
\begin{enumerate}
    \item Definir gene-protein-reaction (GPR) associations
    \item Remover reação(ões) do gene deletado
    \item Resolver FBA com nova rede
    \item Comparar crescimento: $\mu_{KO} / \mu_{WT}$
\end{enumerate}
\end{block}

\textbf{Classificação:}
\begin{itemize}
    \item \textbf{Essencial:} $\mu_{KO} < \epsilon$ (crescimento nulo/mínimo)
    \item \textbf{Não-essencial:} $\mu_{KO} \approx \mu_{WT}$
    \item \textbf{Parcialmente essencial:} $0 < \mu_{KO} < \mu_{WT}$
\end{itemize}

\vspace{0.3cm}

\importante{
FBA prediz corretamente $\sim$85-90\% dos genes essenciais em E. coli em glicose aeróbica
}
\end{frame}

\begin{frame}[fragile]{Implementação: FBA com COBRApy}
\begin{lstlisting}[language=Python, caption=FBA básico com COBRApy]
import cobra
from cobra.io import load_model

# Carregar modelo de E. coli
model = load_model("iJO1366")

print(f"Modelo: {model.id}")
print(f"Reacoes: {len(model.reactions)}")
print(f"Metabolitos: {len(model.metabolites)}")
print(f"Genes: {len(model.genes)}")

# Definir meio de cultura (glicose aerobica)
model.medium = {
    "EX_glc__D_e": 10.0,    # glicose
    "EX_o2_e": 20.0,        # oxigenio
}

# Resolver FBA
solution = model.optimize()

print(f"\nStatus: {solution.status}")
print(f"Taxa de crescimento: {solution.objective_value:.3f} 1/h")

# Fluxos principais
print("\nFluxos principais:")
for rxn_id in ["EX_glc__D_e", "EX_o2_e", "EX_co2_e",
                "EX_h2o_e", "BIOMASS_Ec_iJO1366_core_53p95M"]:
    flux = solution.fluxes[rxn_id]
    print(f"  {rxn_id}: {flux:.2f}")
\end{lstlisting}
\end{frame}

\begin{frame}[fragile]{Implementação: FVA e Gene Knockout}
\begin{lstlisting}[language=Python, caption=FVA e análise de essencialidade, basicstyle=\ttfamily\footnotesize]
from cobra.flux_analysis import flux_variability_analysis
from cobra.flux_analysis import single_gene_deletion

# Flux Variability Analysis
fva_result = flux_variability_analysis(
    model,
    fraction_of_optimum=0.9
)

# Exibir FVA para reacoes centrais
print("FVA (90% do otimo):")
for rxn in ["PGI", "PFK", "FBA", "PYK"]:  # glicolise
    min_flux = fva_result.loc[rxn, "minimum"]
    max_flux = fva_result.loc[rxn, "maximum"]
    print(f"  {rxn}: [{min_flux:.2f}, {max_flux:.2f}]")

# Analise de essencialidade genica
deletion_results = single_gene_deletion(model)

# Classificar genes
essential = deletion_results[
    deletion_results["growth"] < 0.01
]
print(f"\nGenes essenciais: {len(essential)}/{len(model.genes)}")
print("Exemplos:")
for gene in essential.head(5)["ids"]:
    print(f"  {gene}")
\end{lstlisting}
\end{frame}

\begin{frame}[fragile]{Implementação: Visualização de Fluxos}
\begin{lstlisting}[language=Python, caption=Visualizar distribuição de fluxos]
import matplotlib.pyplot as plt
import pandas as pd

# Obter fluxos da solucao
fluxes = solution.fluxes

# Selecionar subsistema (ex: glicolise)
glycolysis = model.groups.get_by_id("Glycolysis")
glyc_rxns = [rxn.id for rxn in glycolysis.members]
glyc_fluxes = fluxes[glyc_rxns]

# Plotar
fig, ax = plt.subplots(figsize=(10, 6))
glyc_fluxes.sort_values().plot(kind='barh', ax=ax)
ax.set_xlabel('Fluxo (mmol/gDW/h)')
ax.set_title('Fluxos na Glicolise')
ax.grid(True, alpha=0.3)
plt.tight_layout()

# Heatmap de fluxos em multiplas condicoes
conditions = ["aerobic", "anaerobic", "minimal"]
flux_matrix = pd.DataFrame()

for condition in conditions:
    # Configurar condicao e resolver
    # ... (codigo de configuracao)
    sol = model.optimize()
    flux_matrix[condition] = sol.fluxes[glyc_rxns]

# Heatmap
import seaborn as sns
plt.figure(figsize=(8, 10))
sns.heatmap(flux_matrix, cmap='RdBu_r', center=0)
plt.title('Fluxos em Diferentes Condicoes')
\end{lstlisting}
\end{frame}

% ============================================================================
% SEÇÃO 5: MODOS ELEMENTARES E VIAS EXTREMAS
% ============================================================================
\section{Modos Elementares e Vias Extremas}

\begin{frame}{Modos Elementares de Fluxo (EFMs)}
\definicao{Elementary Flux Modes}{
Conjuntos mínimos de reações que podem operar em estado estacionário, representando vias metabólicas funcionalmente independentes.
}

\vspace{0.3cm}

\begin{block}{Propriedades}
\begin{itemize}
    \item \textbf{Estado estacionário:} $\mathbf{S} \cdot \vect{v} = \vect{0}$
    \item \textbf{Irreversibilidade:} respeitam termodinâmica
    \item \textbf{Minimalidade:} remoção de qualquer reação quebra o estado estacionário
    \item \textbf{Geneticamente independentes:} cada EFM é único
\end{itemize}
\end{block}

\importante{
EFMs representam todas as rotas metabólicas possíveis em uma rede. Qualquer distribuição de fluxo pode ser expressa como combinação linear de EFMs.
}
\end{frame}

\begin{frame}{Vias Extremas (Extreme Pathways)}
\definicao{Extreme Pathways}{
Subconjunto dos EFMs que forma uma base convexa do espaço de fluxos viáveis. São os vetores geradores do cone de fluxos.
}

\vspace{0.3cm}

\begin{block}{Diferenças: EFMs vs Extreme Pathways}
\begin{center}
\begin{tabular}{lcc}
\toprule
\textbf{Característica} & \textbf{EFMs} & \textbf{ExPas} \\
\midrule
Reversibilidade & Não consideram & Decompõem em irreversível \\
Número & Maior & Menor \\
Base matemática & Vetores suporte & Geradores do cone \\
Interpretação & Vias biológicas & Base convexa \\
\bottomrule
\end{tabular}
\end{center}
\end{block}

\textbf{Relação:}
\begin{itemize}
    \item Extreme pathways $\subseteq$ Elementary modes
    \item ExPas formam base minimal, EFMs incluem combinações
\end{itemize}
\end{frame}

\begin{frame}{Cálculo de EFMs}
\begin{block}{Algoritmos}
\begin{enumerate}
    \item \textbf{Nullspace approach:} Baseado no espaço nulo de $\mathbf{S}$
    \item \textbf{Double description method:} Construção iterativa
    \item \textbf{Bit pattern trees:} Eficiente para redes grandes
\end{enumerate}
\end{block}

\vspace{0.3cm}

\begin{alertblock}{Desafio Computacional}
\begin{itemize}
    \item Número de EFMs cresce exponencialmente com tamanho da rede
    \item Redes pequenas: milhares de EFMs
    \item Redes médias: milhões a bilhões de EFMs
    \item Genoma completo: computacionalmente intratável
\end{itemize}

Limitado a sub-redes ($\sim$50-100 reações)
\end{alertblock}
\end{frame}

\begin{frame}{Interpretação Biológica de EFMs}
\begin{exampleblock}{Aplicações}
\begin{itemize}
    \item \textbf{Identificação de vias:} descobrir rotas metabólicas
    \item \textbf{Análise de robustez:} redundância funcional
    \item \textbf{Engenharia metabólica:} alvos de modificação
    \item \textbf{Análise estrutural:} capacidades metabólicas
    \item \textbf{Otimalidade:} comparar estratégias metabólicas
\end{itemize}
\end{exampleblock}

\vspace{0.3cm}

\exemplo{Metabolismo Central}{
Em E. coli, $\sim$500 EFMs conectam glicose a biomassa, representando estratégias alternativas de:
\begin{itemize}
    \item Produção de ATP (respiração vs fermentação)
    \item Geração de poder redutor (NADPH via PP pathway vs TCA)
    \item Síntese de precursores biossintéticos
\end{itemize}
}
\end{frame}

\begin{frame}{Exemplo: Rede Simples}
\begin{columns}
\column{0.45\textwidth}
\textbf{Rede:}
\begin{center}
\begin{tikzpicture}[scale=0.8, every node/.style={font=\small}]
    \node[draw, circle] (A) at (0,0) {A};
    \node[draw, circle] (B) at (2,1) {B};
    \node[draw, circle] (C) at (2,-1) {C};
    \node[draw, circle] (D) at (4,0) {D};

    \draw[->, thick] (-1,0) -- (A) node[midway, above, font=\tiny] {$v_1$};
    \draw[->, thick] (A) -- (B) node[midway, above left, font=\tiny] {$v_2$};
    \draw[->, thick] (A) -- (C) node[midway, below left, font=\tiny] {$v_3$};
    \draw[->, thick] (B) -- (D) node[midway, above right, font=\tiny] {$v_4$};
    \draw[->, thick] (C) -- (D) node[midway, below right, font=\tiny] {$v_5$};
    \draw[->, thick] (D) -- (5,0) node[midway, above, font=\tiny] {$v_6$};
\end{tikzpicture}
\end{center}

\textbf{Reações:}
\begin{align*}
v_1&: \rightarrow A\\
v_2&: A \rightarrow B\\
v_3&: A \rightarrow C\\
v_4&: B \rightarrow D\\
v_5&: C \rightarrow D\\
v_6&: D \rightarrow
\end{align*}

\column{0.55\textwidth}
\textbf{Matriz $\mathbf{S}$:}
\begin{equation*}
\begin{bmatrix}
1 & -1 & -1 & 0 & 0 & 0\\
0 & 1 & 0 & -1 & 0 & 0\\
0 & 0 & 1 & 0 & -1 & 0\\
0 & 0 & 0 & 1 & 1 & -1
\end{bmatrix}
\end{equation*}

\textbf{Elementary Flux Modes:}
\begin{align*}
\text{EFM}_1&: v_1 = v_2 = v_4 = v_6 = 1, \, v_3 = v_5 = 0\\
&\text{Via: } A \rightarrow B \rightarrow D\\
\\
\text{EFM}_2&: v_1 = v_3 = v_5 = v_6 = 1, \, v_2 = v_4 = 0\\
&\text{Via: } A \rightarrow C \rightarrow D
\end{align*}

\textbf{Interpretação:}
Duas rotas independentes de A para D
\end{columns}
\end{frame}

\begin{frame}[fragile]{Implementação: Cálculo de EFMs (Conceitual)}
\begin{lstlisting}[language=Python, caption=EFMs com efmtool (interface conceitual), basicstyle=\ttfamily\footnotesize]
# Nota: efmtool e uma ferramenta Java
# Aqui mostramos interface Python conceitual

import numpy as np
from scipy.linalg import null_space

# Matriz estequiometrica
S = np.array([
    [ 1, -1, -1,  0,  0,  0],  # A
    [ 0,  1,  0, -1,  0,  0],  # B
    [ 0,  0,  1,  0, -1,  0],  # C
    [ 0,  0,  0,  1,  1, -1]   # D
])

# Calcular espaco nulo (base de solucoes)
null_S = null_space(S)
print("Espaco nulo de S:")
print(null_S)

# EFMs sao combinacoes desses vetores
# que respeitam irreversibilidade (v >= 0)

# Para calculo real de EFMs, usar ferramentas:
# - efmtool (Java): https://csb.ethz.ch/tools/efmtool
# - CellNetAnalyzer (MATLAB)
# - Python wrapper: efmtool-py

# Exemplo de uso conceitual:
# from efmtool import calculate_efms
# efms = calculate_efms(S, reversibilities)
# print(f"Numero de EFMs: {len(efms)}")
# for i, efm in enumerate(efms):
#     print(f"EFM {i+1}: {efm}")
\end{lstlisting}
\end{frame}

% ============================================================================
% SEÇÃO 6: METABOLÔMICA
% ============================================================================
\section{Metabolômica}

\begin{frame}{O que é Metabolômica?}
\definicao{Metabolômica}{
Estudo abrangente do conjunto completo de metabólitos (metaboloma) em uma amostra biológica em um dado momento.
}

\vspace{0.3cm}

\begin{columns}
\column{0.5\textwidth}
\textbf{Características:}
\begin{itemize}
    \item Snapshot do estado metabólico
    \item Reflete fenótipo bioquímico
    \item Integra genoma, transcriptoma, proteoma
    \item Sensível a perturbações ambientais
\end{itemize}

\column{0.5\textwidth}
\textbf{Tipos de Metabólitos:}
\begin{itemize}
    \item \textbf{Primários:} essenciais (glicose, ATP)
    \item \textbf{Secundários:} não-essenciais
    \item \textbf{Xenobióticos:} externos (drogas)
    \item Faixa: $\sim$100 Da a $\sim$1500 Da
\end{itemize}
\end{columns}

\vspace{0.3cm}

\importante{
Metaboloma humano: $\sim$114,000 metabólitos conhecidos (HMDB), $\sim$2,500 endógenos detectáveis
}
\end{frame}

\begin{frame}{Plataformas Analíticas}
\begin{block}{LC-MS: Liquid Chromatography-Mass Spectrometry}
\begin{itemize}
    \item \textbf{Princípio:} separação por cromatografia + detecção por massa
    \item \textbf{Vantagens:} alta sensibilidade, ampla cobertura
    \item \textbf{Aplicações:} metabólitos polares e apolares
    \item \textbf{Métodos:} HILIC, RP, tandem MS (MS/MS)
\end{itemize}
\end{block}

\begin{block}{GC-MS: Gas Chromatography-Mass Spectrometry}
\begin{itemize}
    \item \textbf{Princípio:} volatilização + separação + MS
    \item \textbf{Vantagens:} reprodutibilidade, bibliotecas extensas
    \item \textbf{Desvantagens:} requer derivatização
    \item \textbf{Aplicações:} metabólitos voláteis e derivatizáveis
\end{itemize}
\end{block}

\begin{block}{NMR: Nuclear Magnetic Resonance}
\begin{itemize}
    \item \textbf{Princípio:} ressonância magnética nuclear
    \item \textbf{Vantagens:} não-destrutivo, quantitativo, identificação estrutural
    \item \textbf{Desvantagens:} menor sensibilidade
\end{itemize}
\end{block}
\end{frame}

\begin{frame}{Targeted vs Untargeted Metabolomics}
\begin{columns}
\column{0.5\textwidth}
\begin{block}{Targeted}
\textbf{Foco:} conjunto predefinido de metabólitos

\vspace{0.2cm}

\textbf{Vantagens:}
\begin{itemize}
    \item Alta precisão e acurácia
    \item Quantificação absoluta
    \item Validação com padrões
\end{itemize}

\vspace{0.2cm}

\textbf{Desvantagens:}
\begin{itemize}
    \item Cobertura limitada
    \item Viés de hipótese
\end{itemize}

\vspace{0.2cm}

\textbf{Aplicações:}
\begin{itemize}
    \item Validação de biomarcadores
    \item Monitoramento de vias
\end{itemize}
\end{block}

\column{0.5\textwidth}
\begin{block}{Untargeted}
\textbf{Foco:} detecção abrangente

\vspace{0.2cm}

\textbf{Vantagens:}
\begin{itemize}
    \item Descoberta sem viés
    \item Ampla cobertura
    \item Geração de hipóteses
\end{itemize}

\vspace{0.2cm}

\textbf{Desvantagens:}
\begin{itemize}
    \item Quantificação relativa
    \item Desafios de identificação
    \item Grande volume de dados
\end{itemize}

\vspace{0.2cm}

\textbf{Aplicações:}
\begin{itemize}
    \item Estudos exploratórios
    \item Descoberta de biomarcadores
\end{itemize}
\end{block}
\end{columns}
\end{frame}

\begin{frame}{Pipeline de Análise Metabolômica}
\begin{enumerate}
    \item \textbf{Preparação de amostras}
    \begin{itemize}
        \item Quenching metabólico (parar reações)
        \item Extração de metabólitos
        \item Normalização (por massa, proteína, etc.)
    \end{itemize}

    \item \textbf{Aquisição de dados}
    \begin{itemize}
        \item Análise por LC-MS, GC-MS ou NMR
        \item Controles de qualidade (QC)
    \end{itemize}

    \item \textbf{Processamento de dados}
    \begin{itemize}
        \item Detecção de picos
        \item Alinhamento de tempos de retenção
        \item Normalização e correção de batch
    \end{itemize}

    \item \textbf{Identificação de metabólitos}
    \begin{itemize}
        \item Matching com bibliotecas (m/z, RT)
        \item MS/MS para confirmação estrutural
    \end{itemize}

    \item \textbf{Análise estatística e interpretação biológica}
\end{enumerate}
\end{frame}

\begin{frame}{Identificação de Metabólitos}
\begin{block}{Níveis de Confiança (MSI Standards)}
\begin{enumerate}
    \item \textbf{Nível 1:} Identificado - padrão autêntico + 2 propriedades ortogonais
    \item \textbf{Nível 2:} Putativo - match com biblioteca espectral
    \item \textbf{Nível 3:} Classe putativa - características químicas
    \item \textbf{Nível 4:} Desconhecido - m/z sem identificação
\end{enumerate}
\end{block}

\vspace{0.3cm}

\begin{exampleblock}{Bancos de Dados}
\begin{itemize}
    \item \textbf{HMDB:} Human Metabolome Database (\url{hmdb.ca})
    \item \textbf{METLIN:} biblioteca de MS/MS (\url{metlin.scripps.edu})
    \item \textbf{KEGG Compound:} metabólitos em vias (\url{genome.jp/kegg})
    \item \textbf{ChEBI:} Chemical Entities of Biological Interest
    \item \textbf{PubChem:} NIH database
\end{itemize}
\end{exampleblock}
\end{frame}

\begin{frame}{Integração Multi-Omics}
\begin{center}
\begin{tikzpicture}[node distance=2cm, scale=0.9, transform shape]
    \node[draw, ellipse, fill=roxodna!30] (gen) {Genômica};
    \node[draw, ellipse, fill=laranjacel!30, right=of gen] (trans) {Transcriptômica};
    \node[draw, ellipse, fill=verdeacido!30, below=of gen] (prot) {Proteômica};
    \node[draw, ellipse, fill=azulclaro!30, below=of trans] (met) {Metabolômica};
    \node[draw, rectangle, fill=yellow!30, below right=1cm and -0.5cm of prot] (feno) {Fenótipo};

    \draw[->, thick] (gen) -- (trans);
    \draw[->, thick] (trans) -- (prot);
    \draw[->, thick] (prot) -- (met);
    \draw[->, thick] (met) -- (feno);
    \draw[->, thick, bend left=20] (gen) to (prot);
    \draw[->, thick, bend right=20] (met) to node[right, font=\tiny] {feedback} (trans);
\end{tikzpicture}
\end{center}

\vspace{0.3cm}

\textbf{Abordagens Integrativas:}
\begin{itemize}
    \item \textbf{Mapeamento em vias:} correlacionar níveis de RNA, proteína e metabólito
    \item \textbf{Integração com GEMs:} contextualizar medições metabolômicas
    \item \textbf{Machine learning:} predição de fenótipo multi-omics
    \item \textbf{Network analysis:} identificar módulos regulatórios
\end{itemize}
\end{frame}

\begin{frame}[fragile]{Implementação: Análise de Dados Metabolômicos}
\begin{lstlisting}[language=Python, caption=Processamento básico de dados, basicstyle=\ttfamily\footnotesize]
import pandas as pd
import numpy as np
from sklearn.preprocessing import StandardScaler
from sklearn.decomposition import PCA
import matplotlib.pyplot as plt

# Carregar dados (linhas=amostras, colunas=metabolitos)
data = pd.read_csv("metabolomics_data.csv", index_col=0)
metadata = pd.read_csv("sample_metadata.csv", index_col=0)

print(f"Amostras: {data.shape[0]}, Metabolitos: {data.shape[1]}")

# Pre-processamento
# 1. Remover metabolitos com muitos valores faltantes
threshold = 0.2
data_clean = data.dropna(thresh=len(data) * (1-threshold), axis=1)

# 2. Imputacao de valores faltantes (mediana)
data_imputed = data_clean.fillna(data_clean.median())

# 3. Transformacao logaritmica (dados MS sao log-normais)
data_log = np.log2(data_imputed + 1)

# 4. Normalizacao (Z-score)
scaler = StandardScaler()
data_scaled = pd.DataFrame(
    scaler.fit_transform(data_log),
    index=data_log.index,
    columns=data_log.columns
)

print("Dados processados!")
\end{lstlisting}
\end{frame}

\begin{frame}[fragile]{Implementação: PCA e Visualização}
\begin{lstlisting}[language=Python, caption=Análise de componentes principais, basicstyle=\ttfamily\footnotesize]
# PCA
pca = PCA(n_components=2)
pca_result = pca.fit_transform(data_scaled)

# Variancia explicada
print(f"PC1: {pca.explained_variance_ratio_[0]*100:.1f}%")
print(f"PC2: {pca.explained_variance_ratio_[1]*100:.1f}%")

# Plot PCA score
fig, ax = plt.subplots(figsize=(10, 8))
groups = metadata["group"].unique()
colors = plt.cm.Set1(range(len(groups)))

for group, color in zip(groups, colors):
    mask = metadata["group"] == group
    ax.scatter(pca_result[mask, 0], pca_result[mask, 1],
               c=[color], label=group, s=100, alpha=0.7)

ax.set_xlabel(f'PC1 ({pca.explained_variance_ratio_[0]*100:.1f}%)')
ax.set_ylabel(f'PC2 ({pca.explained_variance_ratio_[1]*100:.1f}%)')
ax.set_title('PCA Score Plot')
ax.legend()
ax.grid(True, alpha=0.3)

# Identificar metabolitos importantes (loadings)
loadings = pd.DataFrame(
    pca.components_.T,
    columns=['PC1', 'PC2'],
    index=data_scaled.columns
)
top_metabolites = loadings['PC1'].abs().nlargest(10)
print("Top metabolitos discriminantes:")
print(top_metabolites)
\end{lstlisting}
\end{frame}

% ============================================================================
% SEÇÃO 7: EXERCÍCIOS
% ============================================================================
\section{Exercícios}

\begin{frame}{Exercícios - Parte 1}
\begin{block}{Questão 1: Cinética Enzimática}
Uma enzima tem $K_M = 5$ mM e $V_{max} = 100$ $\mu$mol/min.
\begin{enumerate}
    \item Calcule a velocidade quando $[S] = 10$ mM
    \item Qual concentração de substrato dá $v = 75\% \cdot V_{max}$?
    \item Se a concentração da enzima dobrar, como mudam $K_M$ e $V_{max}$?
\end{enumerate}
\end{block}

\begin{block}{Questão 2: Balanço de Massa}
Para a rede $A \rightarrow B \rightarrow C$ com $B \rightarrow D$:
\begin{enumerate}
    \item Escreva a matriz estequiométrica $\mathbf{S}$
    \item Escreva as equações de balanço de massa em estado estacionário
    \item Se $v_{A \rightarrow B} = 10$ e $v_{B \rightarrow C} = 6$, calcule $v_{B \rightarrow D}$
\end{enumerate}
\end{block}
\end{frame}

\begin{frame}{Exercícios - Parte 2}
\begin{block}{Questão 3: FBA}
Dado um modelo metabólico simples:
\begin{itemize}
    \item Reação de biomassa requer 40 ATP
    \item Glicose pode gerar ATP via: (a) respiração (38 ATP/glc), (b) fermentação (2 ATP/glc)
    \item Uptake de glicose limitado a 10 mmol/gDW/h
\end{itemize}

\begin{enumerate}
    \item Formule o problema de FBA para maximizar crescimento
    \item Qual é o crescimento máximo se apenas respiração for permitida?
    \item E se apenas fermentação?
    \item Como o modelo escolheria em condições anaeróbicas?
\end{enumerate}
\end{block}
\end{frame}

\begin{frame}{Exercícios - Parte 3}
\begin{alertblock}{Projeto: Análise de GEM com COBRApy}
Escolha um organismo modelo (E. coli, S. cerevisiae, etc.):

\begin{enumerate}
    \item Carregar modelo GEM usando COBRApy
    \item Simular crescimento em meio rico vs mínimo
    \item Realizar análise de gene knockout:
    \begin{itemize}
        \item Identificar genes essenciais
        \item Predizer duplos knockouts sintéticos letais
    \end{itemize}
    \item Realizar FVA para identificar reações flexíveis
    \item Simular produção de um metabólito de interesse:
    \begin{itemize}
        \item Adicionar reação de exportação
        \item Maximizar produção mantendo crescimento mínimo
        \item Identificar modificações genéticas benéficas
    \end{itemize}
\end{enumerate}

Escrever relatório com análises e gráficos!
\end{alertblock}
\end{frame}

% ============================================================================
% SEÇÃO 8: REFERÊNCIAS
% ============================================================================
\section{Referências}

\begin{frame}{Referências Principais - Livros}
\begin{thebibliography}{99}
\bibitem{palsson} Palsson B.O. (2015). \textit{Systems Biology: Constraint-based Reconstruction and Analysis}. Cambridge University Press.

\bibitem{fell} Fell D.A. (1997). \textit{Understanding the Control of Metabolism}. Portland Press.

\bibitem{nielsen} Nielsen J., Jewett M.C. (2008). \textit{Metabolic Engineering}. Springer.

\bibitem{voit} Voit E.O. (2013). \textit{Biochemical Systems Theory: A Review}. ISRN Biomathematics.

\bibitem{stephanopoulos} Stephanopoulos G., Aristidou A.A., Nielsen J. (1998). \textit{Metabolic Engineering: Principles and Methodologies}. Academic Press.
\end{thebibliography}
\end{frame}

\begin{frame}{Referências Principais - Artigos}
\begin{thebibliography}{99}
\bibitem{orth2010} Orth J.D., Thiele I., Palsson B.O. (2010). What is flux balance analysis? \textit{Nature Biotechnology}, 28(3):245-248.

\bibitem{wiechert2001} Wiechert W. (2001). $^{13}$C Metabolic Flux Analysis. \textit{Metabolic Engineering}, 3(3):195-206.

\bibitem{schuster1999} Schuster S. et al. (1999). Detection of elementary flux modes in biochemical networks. \textit{Bioinformatics}, 15(3):251-257.

\bibitem{patti2012} Patti G.J., Yanes O., Siuzdak G. (2012). Metabolomics: the apogee of the omics trilogy. \textit{Nature Reviews Molecular Cell Biology}, 13:263-269.

\bibitem{ebrahim2013} Ebrahim A., Lerman J.A., Palsson B.O., Hyduke D.R. (2013). COBRApy: COnstraints-Based Reconstruction and Analysis for Python. \textit{BMC Systems Biology}, 7:74.
\end{thebibliography}
\end{frame}

\begin{frame}{Ferramentas e Recursos}
\begin{block}{Software}
\begin{itemize}
    \item \textbf{COBRApy:} \url{https://opencobra.github.io/cobrapy}
    \item \textbf{COPASI:} \url{http://copasi.org}
    \item \textbf{CellNetAnalyzer:} \url{https://www2.mpi-magdeburg.mpg.de/projects/cna}
    \item \textbf{efmtool:} \url{https://csb.ethz.ch/tools/efmtool}
    \item \textbf{MetaboAnalyst:} \url{https://www.metaboanalyst.ca}
\end{itemize}
\end{block}

\begin{block}{Bancos de Dados}
\begin{itemize}
    \item \textbf{BiGG Models:} \url{http://bigg.ucsd.edu}
    \item \textbf{KEGG:} \url{https://www.genome.jp/kegg}
    \item \textbf{HMDB:} \url{https://hmdb.ca}
    \item \textbf{METLIN:} \url{https://metlin.scripps.edu}
    \item \textbf{MetaCyc:} \url{https://metacyc.org}
\end{itemize}
\end{block}
\end{frame}

\begin{frame}{Leitura Complementar}
\begin{block}{Tutoriais e Reviews}
\begin{itemize}
    \item Bordbar A. et al. (2014). Constraint-based models predict metabolic and associated cellular functions. \textit{Nature Reviews Genetics}.

    \item Antoniewicz M.R. (2015). Methods and advances in metabolic flux analysis: a mini-review. \textit{Journal of Industrial Microbiology \& Biotechnology}.

    \item Lewis N.E. et al. (2012). Constraining the metabolic genotype-phenotype relationship using a phylogeny of in silico methods. \textit{Nature Reviews Microbiology}.

    \item Zamboni N. et al. (2015). $^{13}$C metabolic flux analysis in complex systems. \textit{Current Opinion in Biotechnology}.
\end{itemize}
\end{block}

\begin{block}{Cursos Online}
\begin{itemize}
    \item \textbf{edX:} Systems Biology (Icahn School of Medicine)
    \item \textbf{Coursera:} Introduction to Computational Biology
    \item \textbf{COBRApy Tutorials:} \url{https://cobrapy.readthedocs.io}
\end{itemize}
\end{block}
\end{frame}

% ============================================================================
% SLIDE FINAL
% ============================================================================
\begin{frame}
\begin{center}
{\Huge Obrigado!}

\vspace{1cm}

{\Large Capítulo 8: Sistemas Metabólicos}

\vspace{1cm}

\textbf{Próximo Capítulo:}\\
Redes de Regulação Gênica

\vspace{1cm}

\textit{Dúvidas? Perguntas?}
\end{center}
\end{frame}

\end{document}
